% Generated from validated Article Draft Result. Do not hand-edit; revise the structured draft instead.
\begingroup\raggedright
\subsection{Agentが「機能」から作業surfaceへ出てくる}
\par\endgroup

4月のagent関連更新を、すでに完成した長時間自律runtimeの到来として読むのは早い。むしろ重要なのは、sandbox、computer use、browsing、memory、workspace automation、orchestrationといった部品が、開発者やチームが日常的に触れるproduct／frameworkへ入り始めたことである。5月以降に実行環境全体の設計が前景化する直前、4月はagentが「modelにtoolを呼ばせる機能」から、作業を預けるsurfaceへ踏み出した段階だった。\autocite{src-1ecdad90bfb4,src-20e34cbfb714,src-c71ef79e3212,src-f9ff221789e0}


\begingroup\raggedright
\subsection{Agents SDK — sandboxとharnessをframework側へ}
\par\endgroup

OpenAIは4月15日、Agents SDKの更新としてnative sandbox executionとmodel-native harnessを公表した。公式説明ではfilesやtoolsを横断するsecureかつlong-runningなagent構築を支えるとしている。ここで注目すべきなのは「長時間動く」という性能主張そのものより、sandboxとharnessがagent frameworkの明示的な構成要素になったことだ。agentを実環境へ接続するほど、model responseだけでなく、どこでcodeを実行し、何にaccessでき、どの状態を保持するかが設計対象になる。\autocite{src-1ecdad90bfb4}


\begingroup\raggedright
\subsection{Codex — coding agentからcomputer workflowへ}
\par\endgroup

4月16日に公表されたCodex appの更新では、computer use、in-app browsing、image generation、memory、pluginsが列挙されている。これらは同じ種類の能力ではない。computer useは画面上の操作、browsingは外部情報への接続、memoryは継続状態、pluginsは拡張interfaceに関わる。したがって「Codexが多機能になった」という要約よりも、codingを起点とするagent surfaceが複数のI/Oと状態管理へ広がったと見る方が、4月から後続月への変化を説明しやすい。\autocite{src-20e34cbfb714}


\begingroup\raggedright
\subsection{Workspace agents — 個人toolからteam workflowへ}
\par\endgroup

4月22日のWorkspace agents in ChatGPTは、Codex-powered agentをcloud上で動かし、複雑なworkflowを複数toolにまたがって自動化するものとして紹介された。一次資料が固定するのはproductのreleaseとstated scopeであり、実際の生産性向上を独立に保証するものではない。それでも、agentの配置先が個人のinteractive sessionだけでなくteam workspaceへ広がったことは、authorization、共有context、運用境界が後に重要になる流れの一部として読める。\autocite{src-c71ef79e3212}


\begingroup\raggedright
\subsection{Symphony — issue trackerをorchestration surfaceにする}
\par\endgroup

4月27日にOpenAIが公開したSymphonyは、Codex orchestrationのopen-source specとして、issue trackerをalways-on agent systemへ接続する構想を示す。ここでも「always-on」はagentが無制限に自律するという意味ではなく、仕事の入口、task queue、実行主体をorchestrationとして結び付ける設計語として読むべきだ。issue単位の仕事をagentへ渡す発想は、単発promptからwork queueへと運用単位が変わる兆候でもある。\autocite{src-f9ff221789e0}


\begingroup\raggedright
\subsection{4月時点で見えるstackの層}
\par\endgroup

4件を並べると、同じagentという言葉の下に異なる層がある。Agents SDKはexecution substrateとharness、Codex appはuser-facing interaction surface、Workspace agentsはteam-level deployment、Symphonyはtask orchestrationを扱う。これらを単一製品のfeature listへ畳まず層として見ると、5月以降にsandbox、権限、retrieval、serving、persistent stateが「実行環境」として統合的に論じられる理由が見えてくる。\autocite{src-1ecdad90bfb4,src-20e34cbfb714,src-c71ef79e3212,src-f9ff221789e0}


\begingroup\raggedright
\subsection{4月と5月を同じ結論にしない}
\par\endgroup

Retrospectiveで注意したいのは、後から見えている成熟像を4月へ持ち込まないことだ。4月の一次資料から確実に言えるのは、sandbox、computer use、memory、workspace、orchestrationといった構成要素が明示的になったことまでである。これらが長時間agentの運用基盤としてどう組み合わさり、権限・retrieval・servingと一体化するかは5月以降の主題として残す。この段階差があるからこそ、4月を「競争軸の分化」、5月を「実行環境化」として連続して読める。\autocite{src-1ecdad90bfb4,src-20e34cbfb714,src-c71ef79e3212,src-f9ff221789e0}


